%% kapa: exam=1 round=23 year=2012 period=2 form=A subject=appraisal_law subject_ko=부동산관계법규 seq=29 number=69
%% source: Q-Net 감정평가사 기출문제 (cst00309) artlSeq=5052580 file=raw/1008/5052580/2012년 제23회 감정평가사 제1차시험 기출문제(2교시 A형).pdf
%% conversion: kordoc 4.24.51 -> tools/gen1.py (auto); review=proofread-sonnet; flags=none
\begin{question}{69}
건축법령상 용어의 정의로서 옳지 않은 것은?
\begin{choices}
  \item “고층건축물”이란 층수가 30층 이상이거나 높이가 100m 이상인 건축물을 말한다.
  \item “지하층”이란 건축물의 바닥이 지표면 아래에 있는 층으로서 바닥에서 지표면까지 평균높이가 해당 층 높이의 2분의 1 이상인 것을 말한다.
  \item “리모델링”이란 건축물의 노후화를 억제하거나 기능 향상 등을 위하여 대수선하거나 일부 증축하는 행위를 말한다.
  \item “건축물의 용도”란 건축물의 종류를 유사한 구조, 이용 목적 및 형태별로 묶어 분류한 것을 말한다.
  \item “건축주”란 건축물의 건축ㆍ대수선ㆍ용도변경, 건축설비의 설치 또는 공작물의 축조에 관한 공사를 발주하거나 현장 관리인을 두어 스스로 그 공사를 하는 자를 말한다.
\end{choices}
\end{question}
\answer{1}
